%% GENERATED by selfunc_v411.py -- do not hand-edit.
\begin{table}
\centering
\scriptsize
\setlength{\tabcolsep}{4pt}
\caption{Selection function of the searched sample against a volume-limited reference census, both selected on the same parallax quality: Gaia~DR3 sources with $\varpi>0$, $\sigma_\varpi/\varpi\le0.1$ and $1/\varpi\le40$\,pc (17\,566 objects, 8594 of them carrying a \texttt{teff\_gspphot} estimate), against the 89 stars with at least one searched window. The searched stars are measured far better than the cut requires --- the worst fractional parallax error among them is 2.5 per cent --- so the two columns are selected alike, and loosening the cut to 0.2 adds no census object at all. 86 of the 89 searched stars are classified: 53 from the Gaia temperature, 25 from the SIMBAD spectral type where that temperature is absent, and 8 identified as white dwarfs; 3 carry no classification in either source. The SIMBAD route matters because \texttt{teff\_gspphot} is missing preferentially for cool stars, which is the population the M~dwarf statement is about. ``Ratio'' is the sample column fraction over the census column fraction; $>1$ is over-representation. \emph{The strongest selection effect here is distance, not spectral type.}}
\label{tab:selfunc}
\begin{tabular}{@{}lrrr@{}}
\toprule
Property & Census & Searched & Ratio \\
\midrule
\multicolumn{4}{@{}l}{Distance} \\
\quad $d\le5$\,pc & 48 (0.3\%) & 12 (13.5\%) & 49.3 \\
\quad 5--10\,pc & 240 (1.4\%) & 13 (14.6\%) & 10.7 \\
\quad 10--20\,pc & 2069 (11.8\%) & 18 (20.2\%) & 1.7 \\
\quad 20--30\,pc & 5391 (30.7\%) & 19 (21.3\%) & 0.7 \\
\quad 30--40\,pc & 9818 (55.9\%) & 27 (30.3\%) & 0.5 \\
\midrule
\multicolumn{4}{@{}l}{$T_{\rm eff}$ class$^{a}$} \\
\quad A & 24 (0.3\%) & 5 (6.4\%) & 23.0 \\
\quad F & 402 (4.7\%) & 9 (11.5\%) & 2.5 \\
\quad G & 860 (10.0\%) & 16 (20.5\%) & 2.0 \\
\quad K & 1400 (16.3\%) & 7 (9.0\%) & 0.6 \\
\quad M & 5908 (68.7\%) & 41 (52.6\%) & 0.8 \\
\midrule
\multicolumn{4}{@{}l}{Other axes} \\
\quad Confirmed planet hosts$^{b}$ & 26 (15.5\%) & 23 (25.8\%) & 1.7 \\
\quad White dwarfs$^{a}$ & -- & 8 (9.0\%) & -- \\
\quad Searched stars / systems & -- & 89 / 82 & -- \\
\bottomrule
\end{tabular}

\smallskip
{\footnotesize $^{a}$Temperature bins, not MK types: A $\ge7500$, F 6000--7500, G 5300--6000, K 3900--5300, M $<3900$\,K, normalised over the 78 classified non-degenerate stars. The 8 white dwarfs are kept apart because a degenerate is not an early-type star and its photometric temperature would place it among them. Gaia \texttt{teff\_gspphot} is unavailable for about half the reference census, preferentially for the coolest and faintest objects, so the census M fraction is a lower limit and every earlier-type ratio in this panel is correspondingly a lower limit; the searched column has no such gap.\\ $^{b}$Both columns are the same positional join, of the 168-entry ALMA-covered work list and of its searched subset, to the NASA Exoplanet Archive \texttt{pscomppars} table accessed 2026 October 4; a confirmed planet is a row of that table. The 23 searched hosts carry 50 confirmed planets between them, and 21 of the 23 survive the deuterium-burning mass limit. The Gaia reference census carries no planet information, so the census column in this row is the work list rather than the 17\,566-object Gaia census.}
\end{table}
